\documentclass[a4paper,11pt]{article}

\usepackage[a4paper,margin=2cm]{geometry}
\usepackage[T1]{fontenc}
\usepackage[utf8]{inputenc}
\usepackage[ngerman,english]{babel}
\usepackage{lmodern}
\usepackage{microtype}
\usepackage[table]{xcolor}
\usepackage{parskip}
\usepackage{enumitem}
\usepackage[most]{tcolorbox}
\usepackage{titlesec}
\usepackage{array}
\usepackage{longtable}
\usepackage[hidelinks]{hyperref}

\definecolor{HeaderBlueDark}{RGB}{70,110,170}
\definecolor{RowBlueLight}{RGB}{235,243,255}
\definecolor{RowBlueDark}{RGB}{220,233,250}
\definecolor{SoftGray}{RGB}{90,90,90}

\setlength{\parindent}{0pt}
\setlist[itemize]{leftmargin=1.4em,itemsep=0.35em,topsep=0.35em}
\linespread{1.06}
\renewcommand{\arraystretch}{1.35}
\setlength{\tabcolsep}{6pt}

\titleformat{\section}[block]
{\Large\bfseries\color{HeaderBlueDark}}
{}{0pt}{}
\titlespacing{\section}{0pt}{1.3em}{0.8em}

\titleformat{\subsection}[block]
{\large\bfseries\color{HeaderBlueDark}}
{}{0pt}{}
\titlespacing{\subsection}{0pt}{1.0em}{0.6em}

\newtcolorbox{exbox}{enhanced,breakable,colback=RowBlueLight,colframe=RowBlueDark,boxrule=0.4pt,arc=1.5mm,left=1.2mm,right=1.2mm,top=1mm,bottom=1mm,title={Beispiele \textnormal{(Examples)}},fonttitle=\bfseries,colbacktitle=RowBlueDark,coltitle=black}

\NewDocumentEnvironment{grammarentry}{mm+b}{%
    \begin{tcolorbox}[enhanced,breakable,colback=white,colframe=HeaderBlueDark,boxrule=0.6pt,arc=1.5mm,left=1.5mm,right=1.5mm,top=1mm,bottom=1mm,colbacktitle=HeaderBlueDark,coltitle=white,fonttitle=\bfseries,title={#1\hfill\normalfont\footnotesize #2}]
    #3
    \end{tcolorbox}
}{}

\newcommand{\GermanText}[1]{\textbf{Deutsch: }#1\par}
\newcommand{\EnglishText}[1]{{\small\color{SoftGray}\textbf{English: }#1\par}}
\newcommand{\ExampleItem}[2]{\item \textbf{#1}\\{\small\color{SoftGray}#2}}

\title{Adverbiale Adjektive\\\large\textnormal{Adverbial Adjectives}}
\date{}

\begin{document}
    \maketitle
    \tableofcontents
    \newpage

    \section{Grundidee \textnormal{(Basic idea)}}

        \begin{grammarentry}{Adverbialer Gebrauch}{Funktion eines Adjektivs}
            \GermanText{Ein \textbf{adverbial gebrauchtes Adjektiv} beschreibt nicht direkt ein Nomen, sondern vor allem, \textbf{wie} eine Handlung oder ein Vorgang geschieht. Es bezieht sich typischerweise auf das Verb.}
            \EnglishText{An \textbf{adverbially used adjective} does not directly describe a noun. Instead, it mainly describes \textbf{how} an action or process happens. It typically relates to the verb.}

            \GermanText{Das wichtigste Merkmal ist: Das Adjektiv bekommt im adverbialen Gebrauch \textbf{keine Adjektivendung}.}
            \EnglishText{The most important feature is: In adverbial use, the adjective takes \textbf{no adjective ending}.}
        \end{grammarentry}

        \begin{exbox}
            \begin{itemize}
                \ExampleItem{Er fährt schnell.}{He drives quickly.}
                \ExampleItem{Sie spricht leise.}{She speaks quietly.}
                \ExampleItem{Die Kinder arbeiten konzentriert.}{The children work in a focused way.}
            \end{itemize}
        \end{exbox}

    \section{Form: keine Deklination \textnormal{(Form: no declension)}}

        \begin{grammarentry}{Unflektierte Form}{keine Endung}
            \GermanText{Adverbial gebrauchte Adjektive werden \textbf{nicht nach Genus, Numerus oder Kasus dekliniert}. Die Form bleibt unverändert.}
            \EnglishText{Adjectives used adverbially are \textbf{not declined according to gender, number, or case}. Their form remains unchanged.}

            \GermanText{Vergleiche: \textbf{ein schnelles Auto} (attributiv, mit Endung) - \textbf{Das Auto fährt schnell} (adverbial, ohne Endung).}
            \EnglishText{Compare: \textbf{ein schnelles Auto} (attributive, with an ending) - \textbf{Das Auto fährt schnell} (adverbial, without an ending).}
        \end{grammarentry}

        \rowcolors{2}{RowBlueLight}{RowBlueDark}
        \begin{longtable}{>{\bfseries}p{3.5cm}|p{5.4cm}|p{5.4cm}}
            \rowcolor{HeaderBlueDark}
            \color{white}\textbf{Gebrauch / Use} &
            \color{white}\textbf{Deutsch / German} &
            \color{white}\textbf{English} \\
            Attributiv / Attributive & ein \textbf{schnelles} Auto & a fast car \\
            Prädikativ / Predicative & Das Auto ist \textbf{schnell}. & The car is fast. \\
            Adverbial / Adverbial & Das Auto fährt \textbf{schnell}. & The car drives quickly. \\
        \end{longtable}

    \section{Woran erkennt man den adverbialen Gebrauch? \textnormal{(How do you recognize adverbial use?)}}

        \begin{grammarentry}{Erkennungsfrage}{meistens: Wie?}
            \GermanText{Ein adverbial gebrauchtes Adjektiv beantwortet sehr häufig die Frage \textbf{Wie?}.}
            \EnglishText{An adjective used adverbially very often answers the question \textbf{How?}.}

            \GermanText{Beispiel: \glqq Wie fährt er?\grqq{} - \glqq Er fährt \textbf{vorsichtig}.\grqq}
            \EnglishText{Example: ``How does he drive?'' - ``He drives \textbf{carefully}.''}

            \GermanText{Diese Frageprobe ist sehr hilfreich, aber nicht jede adverbiale Bestimmung muss ein Adjektiv sein. Auch echte Adverbien oder Präpositionalgruppen können dieselbe Funktion übernehmen.}
            \EnglishText{This question test is very useful, but not every adverbial modifier has to be an adjective. True adverbs or prepositional phrases can perform the same function.}
        \end{grammarentry}

    \section{Adjektiv oder Adverb? \textnormal{(Adjective or adverb?)}}

        \begin{grammarentry}{Wichtige Besonderheit im Deutschen}{gleiche Form}
            \GermanText{Im Deutschen braucht man für den adverbialen Gebrauch normalerweise \textbf{keine besondere Adverbform}. Das normale Adjektiv wird unverändert verwendet.}
            \EnglishText{In German, adverbial use normally does \textbf{not require a special adverb form}. The ordinary adjective is used unchanged.}

            \GermanText{Deshalb entspricht deutsches \textbf{schnell} je nach Satz sowohl englischem \textbf{fast} als auch einer englischen adverbialen Bedeutung wie \textbf{quickly}.}
            \EnglishText{Therefore, German \textbf{schnell} can correspond, depending on the sentence, both to English \textbf{fast} and to an English adverbial meaning such as \textbf{quickly}.}
        \end{grammarentry}

        \begin{exbox}
            \begin{itemize}
                \ExampleItem{Er ist ruhig.}{He is calm. - predicative use}
                \ExampleItem{Er spricht ruhig.}{He speaks calmly. - adverbial use}
                \ExampleItem{Sie erklärt die Aufgabe klar.}{She explains the task clearly.}
            \end{itemize}
        \end{exbox}

    \section{Steigerung \textnormal{(Comparison)}}

        \begin{grammarentry}{Komparativ und Superlativ}{auch adverbial möglich}
            \GermanText{Wenn ein Adjektiv steigerbar ist, kann es auch im adverbialen Gebrauch im \textbf{Komparativ} und \textbf{Superlativ} stehen.}
            \EnglishText{If an adjective is gradable, it can also be used adverbially in the \textbf{comparative} and \textbf{superlative}.}

            \GermanText{Komparativ: Grundform + \textbf{-er}, zum Beispiel \textbf{schneller}. Es gibt weiterhin keine zusätzliche Adjektivendung.}
            \EnglishText{Comparative: base form + \textbf{-er}, for example \textbf{schneller}. There is still no additional adjective ending.}

            \GermanText{Beim adverbialen Superlativ verwendet man normalerweise die Form \textbf{am + -sten/-esten}: \textbf{am schnellsten}, \textbf{am besten}, \textbf{am vorsichtigsten}.}
            \EnglishText{For the adverbial superlative, German normally uses \textbf{am + -sten/-esten}: \textbf{am schnellsten}, \textbf{am besten}, \textbf{am vorsichtigsten}.}
        \end{grammarentry}

        \begin{exbox}
            \begin{itemize}
                \ExampleItem{Anna läuft schnell.}{Anna runs fast.}
                \ExampleItem{Lukas läuft schneller.}{Lukas runs faster.}
                \ExampleItem{Mira läuft am schnellsten.}{Mira runs the fastest.}
            \end{itemize}
        \end{exbox}

    \section{Stellung im Satz \textnormal{(Position in the sentence)}}

        \begin{grammarentry}{Keine einzige feste Position}{abhängig vom Satz}
            \GermanText{Ein adverbial gebrauchtes Adjektiv hat nicht eine einzige feste Position. Es steht innerhalb des Satzes dort, wo die adverbiale Information sinnvoll eingeordnet wird.}
            \EnglishText{An adverbially used adjective does not have one single fixed position. It appears within the sentence where the adverbial information is naturally placed.}

            \GermanText{Bei einfachen Sätzen steht es oft im Mittelfeld oder nahe beim Verb beziehungsweise bei dem Teil, den es näher bestimmt.}
            \EnglishText{In simple sentences, it often appears in the middle field or close to the verb or the element it modifies more precisely.}

            \GermanText{Die allgemeinen Regeln der deutschen Wortstellung, besonders die Reihenfolge mehrerer adverbialer Angaben, gelten weiterhin.}
            \EnglishText{The general rules of German word order, especially the ordering of several adverbial elements, still apply.}
        \end{grammarentry}

        \begin{exbox}
            \begin{itemize}
                \ExampleItem{Er hat die Frage ruhig beantwortet.}{He answered the question calmly.}
                \ExampleItem{Sie arbeitet heute besonders konzentriert.}{She is working especially attentively today.}
                \ExampleItem{Vorsichtig öffnete er die Tür.}{He carefully opened the door.}
            \end{itemize}
        \end{exbox}

    \section{Erweiterung und Gradangaben \textnormal{(Modification and degree)}}

        \begin{grammarentry}{Sehr, besonders, ziemlich, zu}{Grad des Adjektivs}
            \GermanText{Ein adverbial gebrauchtes Adjektiv kann durch Gradwörter näher bestimmt werden, zum Beispiel \textbf{sehr}, \textbf{besonders}, \textbf{ziemlich}, \textbf{extrem}, \textbf{zu} oder \textbf{genug}.}
            \EnglishText{An adjective used adverbially can be modified by degree words such as \textbf{very}, \textbf{especially}, \textbf{quite}, \textbf{extremely}, \textbf{too}, or \textbf{enough}.}

            \GermanText{Auch dann bleibt das Adjektiv ohne Deklinationsendung.}
            \EnglishText{Even then, the adjective remains without a declensional ending.}
        \end{grammarentry}

        \begin{exbox}
            \begin{itemize}
                \ExampleItem{Sie spricht sehr deutlich.}{She speaks very clearly.}
                \ExampleItem{Er fährt zu schnell.}{He drives too fast.}
                \ExampleItem{Du hast die Aufgabe gut genug erklärt.}{You explained the task well enough.}
            \end{itemize}
        \end{exbox}

    \section{Partizipien im adverbialen Gebrauch \textnormal{(Participles in adverbial use)}}

        \begin{grammarentry}{Partizip I und Partizip II}{adjektivischer Gebrauch}
            \GermanText{Auch Partizipien können wie Adjektive adverbial gebraucht werden. Dann stehen sie ebenfalls ohne Adjektivendung.}
            \EnglishText{Participles can also be used adverbially like adjectives. In that case, they also appear without an adjective ending.}

            \GermanText{Partizip I: \textbf{lächelnd}, \textbf{schweigend}. Partizip II: zum Beispiel \textbf{konzentriert}, \textbf{genervt}, wenn die Form eine Eigenschaft oder einen Zustand ausdrückt.}
            \EnglishText{Present participle: \textbf{lächelnd} (smiling), \textbf{schweigend} (silently). Past participle: for example \textbf{konzentriert} (focused), \textbf{genervt} (annoyed), when the form expresses a property or state.}
        \end{grammarentry}

        \begin{exbox}
            \begin{itemize}
                \ExampleItem{Sie kam lächelnd ins Zimmer.}{She came into the room smiling.}
                \ExampleItem{Er hörte schweigend zu.}{He listened silently.}
                \ExampleItem{Sie arbeitet konzentriert weiter.}{She continues working in a focused way.}
            \end{itemize}
        \end{exbox}

    \section{Abgrenzung zum attributiven Adjektiv \textnormal{(Difference from attributive adjectives)}}

        \begin{grammarentry}{Attributiv = direkt beim Nomen}{mit Deklination}
            \GermanText{Ein attributives Adjektiv steht direkt bei einem Nomen und bekommt eine Endung: \textbf{ein schneller Zug}, \textbf{mit einem schnellen Zug}.}
            \EnglishText{An attributive adjective stands directly with a noun and receives an ending: \textbf{ein schneller Zug} (a fast train), \textbf{mit einem schnellen Zug} (with a fast train).}

            \GermanText{Ein adverbiales Adjektiv beschreibt dagegen die Handlung oder den Vorgang: \textbf{Der Zug fährt schnell}.}
            \EnglishText{An adverbial adjective, by contrast, describes the action or process: \textbf{Der Zug fährt schnell} (The train travels fast).}
        \end{grammarentry}

    \section{Abgrenzung zum prädikativen Adjektiv \textnormal{(Difference from predicative adjectives)}}

        \begin{grammarentry}{Beide ohne Endung}{aber unterschiedliche Funktion}
            \GermanText{Sowohl prädikative als auch adverbiale Adjektive stehen normalerweise \textbf{ohne Adjektivendung}. Der Unterschied liegt deshalb nicht in der Form, sondern in der Funktion.}
            \EnglishText{Both predicative and adverbial adjectives normally appear \textbf{without an adjective ending}. The difference is therefore not in the form but in the function.}

            \GermanText{Prädikativ: Das Adjektiv beschreibt das Subjekt oder ein anderes Satzglied über ein verbindendes Verb: \textbf{Das Auto ist schnell}.}
            \EnglishText{Predicative: The adjective describes the subject or another sentence element through a linking verb: \textbf{Das Auto ist schnell} (The car is fast).}

            \GermanText{Adverbial: Das Adjektiv beschreibt, wie eine Handlung geschieht: \textbf{Das Auto fährt schnell}.}
            \EnglishText{Adverbial: The adjective describes how an action happens: \textbf{Das Auto fährt schnell} (The car drives fast).}
        \end{grammarentry}

    \section{Typische Fehler \textnormal{(Typical mistakes)}}

        \begin{grammarentry}{Fehler 1}{unnötige Adjektivendung}
            \GermanText{Falsch: \textbf{Er fährt schnelles.} Richtig: \textbf{Er fährt schnell.}}
            \EnglishText{Wrong: \textbf{Er fährt schnelles.} Correct: \textbf{Er fährt schnell.}}
        \end{grammarentry}

        \begin{grammarentry}{Fehler 2}{englische -ly-Logik übertragen}
            \GermanText{Im Deutschen bildet man normalerweise kein neues Adverb wie im Englischen mit \textit{-ly}. Man verwendet die unveränderte Adjektivform: \textbf{ruhig}, \textbf{klar}, \textbf{schnell}.}
            \EnglishText{German normally does not form a new adverb as English does with \textit{-ly}. It uses the unchanged adjective form: \textbf{ruhig}, \textbf{klar}, \textbf{schnell}.}
        \end{grammarentry}

        \begin{grammarentry}{Fehler 3}{adverbial und prädikativ verwechseln}
            \GermanText{\textbf{Er ist vorsichtig} ist prädikativ. \textbf{Er fährt vorsichtig} ist adverbial. In beiden Fällen steht \textbf{vorsichtig} ohne Endung, aber die grammatische Funktion ist verschieden.}
            \EnglishText{\textbf{Er ist vorsichtig} is predicative. \textbf{Er fährt vorsichtig} is adverbial. In both cases, \textbf{vorsichtig} has no ending, but the grammatical function is different.}
        \end{grammarentry}

    \section{Schnelltest \textnormal{(Quick test)}}

        \begin{grammarentry}{Drei Fragen}{Erkennung}
            \GermanText{1. Steht das Adjektiv direkt vor einem Nomen und beschreibt dieses? Dann ist es meistens \textbf{attributiv} und wird dekliniert.}
            \EnglishText{1. Does the adjective stand directly before a noun and describe it? Then it is usually \textbf{attributive} and is declined.}

            \GermanText{2. Beschreibt das Adjektiv über \textbf{sein, werden, bleiben} oder ein ähnliches Verb einen Zustand oder eine Eigenschaft des Subjekts? Dann ist es \textbf{prädikativ}.}
            \EnglishText{2. Does the adjective describe a state or property of the subject through \textbf{sein, werden, bleiben} or a similar verb? Then it is \textbf{predicative}.}

            \GermanText{3. Beschreibt das Adjektiv vor allem, \textbf{wie} eine Handlung geschieht? Dann ist es \textbf{adverbial}.}
            \EnglishText{3. Does the adjective mainly describe \textbf{how} an action happens? Then it is \textbf{adverbial}.}
        \end{grammarentry}

    \section{Zusammenfassung \textnormal{(Summary)}}

        \begin{grammarentry}{Merksätze}{kurz und wichtig}
            \GermanText{Adverbial gebrauchte Adjektive beschreiben typischerweise eine Handlung oder einen Vorgang und beantworten oft die Frage \textbf{Wie?}.}
            \EnglishText{Adjectives used adverbially typically describe an action or process and often answer the question \textbf{How?}.}

            \GermanText{Sie haben \textbf{keine Deklinationsendung}: \textbf{schnell, ruhig, deutlich, vorsichtig}.}
            \EnglishText{They have \textbf{no declensional ending}: \textbf{schnell, ruhig, deutlich, vorsichtig}.}

            \GermanText{Komparativ und Superlativ sind möglich: \textbf{schneller}, \textbf{am schnellsten}.}
            \EnglishText{Comparative and superlative forms are possible: \textbf{schneller} (faster), \textbf{am schnellsten} (the fastest).}

            \GermanText{Der wichtigste Gegensatz lautet: \textbf{ein schnelles Auto} (attributiv) - \textbf{Das Auto ist schnell} (prädikativ) - \textbf{Das Auto fährt schnell} (adverbial).}
            \EnglishText{The key contrast is: \textbf{ein schnelles Auto} (attributive) - \textbf{Das Auto ist schnell} (predicative) - \textbf{Das Auto fährt schnell} (adverbial).}
        \end{grammarentry}

\end{document}
